
\subsubsection{1.1 Background}

Reinforcement learning from execution feedback has emerged as a paradigm for improving code generation in language models. CodeRL introduced using program execution outcomes as reward signals. RLTF extended this approach by combining coarse (pass/fail) and fine-grained (token-level) feedback, demonstrating that multi-granularity feedback outperforms single-signal approaches. VeRPO identified that aggregation strategy matters---cardinality bias in dense rewards can degrade performance. These findings establish that feedback granularity and aggregation are critical design choices.

\subsubsection{1.2 The Problem}

Existing methods apply fine-grained feedback uniformly regardless of whether error localization is reliable. When a Python program raises a SyntaxError, the traceback points to the exact buggy token. When it raises a RuntimeError, the traceback often points to a symptom line far from the root cause. Applying token-level penalties at such misleading locations injects noise into gradient updates: the model receives credit assignment signals at wrong tokens.

\subsubsection{1.3 Key Insight}

Python exception types partition into two categories with different localization reliability. In our experiments using synthetic code samples with known bug locations:

\begin{itemize}
\item \textbf{U\_line errors} (SyntaxError, IndentationError, NameError, TypeError, AttributeError, KeyError, IndexError, ValueError, ZeroDivisionError) achieved 100\% traceback accuracy---the reported line matched the actual bug location in all 250 samples tested.
\item \textbf{U\_ignore errors} (RuntimeError, RecursionError, MemoryError, TimeoutError, AssertionError) achieved 20\% accuracy across 250 samples---tracebacks predominantly pointed to symptom locations.
\end{itemize}

This 80-percentage-point gap (chi-square p = 9.57 $\times 10^{-74})$ suggests that conditioning feedback strategy on error type could reduce gradient noise.

\subsubsection{1.4 Our Approach}

We propose error-type-gated fine-grained feedback: apply token-level penalties only for errors where localization is reliable (U\_line category), while falling back to coarse-only feedback for unreliably-localized errors (U\_ignore category). This modification preserves the benefits of fine-grained credit assignment where it is accurate while avoiding noise injection where it would be misleading.

\subsubsection{1.5 Contributions}

Through a systematic verification pipeline, we establish the causal mechanism underlying this approach:

\begin{enumerate}
\item Fine-grained penalties concentrate gradients at traceback locations (16.$11\times$ concentration ratio, p < $10^{-270})$
\item Localization accuracy varies dramatically by error type (100\% vs 20\%, p < $10^{-73})$
\item Unreliable localization causes measurable gradient noise at ground-truth locations (p < $10^{-13})$
\item Gating improves signal-to-noise ratio (+4.61\%, direction confirmed, p = 0.112)
\item In proof-of-concept training, gating enables the model to reach 30\% pass@1 when the baseline does not
\end{enumerate}

All experiments use CodeT5-small (60M parameters) and synthetic code samples for proof-of-concept validation. Generalization to larger models and real-world samples requires further investigation.

